\documentclass[12pt, a4paper]{article}
\usepackage{graphicx} % Required for inserting images
\graphicspath{{images/}{images-desafios/}}
\usepackage[a4paper, total={7.5in, 10in}]{geometry}
\usepackage{float} % pra conseguir usar o \begin{figure}[H]
\usepackage{indentfirst} % pra fazer tab n 1 paragrafos dos capitulos
% package pra tabela
\usepackage{multirow}
\usepackage{array}

\newcommand{\unit}[1]{\ensuremath{\, \mathrm{#1}}} % comando pra tirar do math mode a unidade de medida

\def\tamanhoGrafico{0.8}

\newcommand{\ohm}{\unit{\Omega} }
\newcommand{\kohm}{\unit{k\Omega}}

\title{
	Relatório de Laboratório Digital A - PCS3335
	\\ Experiência 1
}

\author{
	Alunos:
	\\ Rafael Hipólito Rodrigues - N.USP: 14604308
	\\ {{PERSON-2}} - N.USP: {{PERSON-ID-2}}
	\\ \\
	Turma 3 - Bancada A6
	\\ Professor: Bruno de Carvalho Albertini
}

\date{12/03/2025}

\begin{document}
	
	\maketitle
	
	\section*{Execução}
	
	Elaboração da entidade que recebe um número binário de 10 bits das chaves e exibe o número em hexadecimal nos 6 displasys de 7 segmentos de forma palíndroma. 
	
	\subsection*{Diagrama RTL}
	
	\begin{figure}[H]
		\centering
		\includegraphics[width=0.7\linewidth]{images/bolsonaroPresidente2028_vulgo_exp1}
		\caption{Diagrama RTL do VHDL.}
		\label{fig:bolsonaroPresidente2028_vulgo_exp1}
	\end{figure}
	
	\subsection*{Testes}

	\begin{center}
		\begin{tabular}{ |c|c|c| }
			\hline
			Valor dos Switchs(Binário) & Valor esperado(Hexa) & Valor obtido(Hexa) \\
			\hline
			00 0000 0000 & 000000 & 000000 \\
			\hline
			00 0000 0001 & 100001 & 100001 \\
			\hline
			00 0000 0010 & 200002 & 200002 \\
			\hline
			00 0000 1111 & F0000F & F0000F \\
			\hline
			00 0001 0000 & F1001F & F1001F \\
			\hline
			00 1111 1111 & FF00FF & FF00FF \\
			\hline
			01 0000 0000 & FF11FF & FF11FF \\
			\hline
			11 1111 1111 & FF33FF & FF33FF \\
			\hline
		\end{tabular}
	\end{center}
	
	\section*{Desafio}
	
	Elaboração de uma entidade ascii2seg que recebe um número binário de 7 bits das chaves e exibe o correspondente ASCII em um display de 7 segmentos.
	
	\begin{figure}[H]
		\centering
		\includegraphics[width=1\linewidth]{images-desafios/screenshot1.jpg}
		\caption{Diagrama RTL da entidade ascii2seg completo.}
		\label{fig:screenshot1}
	\end{figure}
	
	\begin{figure}[H]
		\centering
		\includegraphics[width=0.8\linewidth]{images-desafios/screenshot2.jpg}
		\caption{Inicio do Diagrama RTL da entidade ascii2seg .}
		\label{fig:screenshot2}
	\end{figure}
	
	\begin{figure}[H]
		\centering
		\includegraphics[width=0.8\linewidth]{images-desafios/screenshot3.jpg}
		\caption{Final do diagrama RTL da entidade ascii2seg.}
		\label{fig:screenshot3}
	\end{figure}
	
	Devido a necessidade de atribuir 96 valores(quantidade de caracteres ASCII utilizados) diferentes a o sinal de saída, o circuito precisou ficou muito longo.
	
	\subsection*{Testes}

	\begin{center}
		\begin{tabular}{ |c|c|c| }
			\hline
			Valor dos Switchs(Binário) & Valor esperado(Hexa) & Valor obtido(Hexa) \\
			\hline
			011 0000 & 0 & 0 \\
			\hline
			011 0001 & 1 & 1 \\
			\hline
			011 1001 & 9 & 9 \\
			\hline
			100 0001 & A & A \\
			\hline
			100 0010 & B & B \\
			\hline
			101 1010 & Z & Z \\
			\hline
		\end{tabular}
	\end{center}

	\subsection*{Observações}
	
	O grupo descobriu durante o experimento que os leds do display de 7 segmentos estavam ligados com a polaridade contraria da qual o grupo esperava. Assim, houve a necessidade de negar(truncar os bits) dos sinais de saída que vão para o display de 7 segmentos. 
	
\end{document}
